\documentclass[12pt]{article}
\usepackage{geometry}
\geometry{margin=1in}
\usepackage{hyperref}

\begin{document}

\noindent
March 2026

\bigskip

\noindent
Professor Mariarosaria Taddeo\\
Editor-in-Chief\\
\textit{Minds and Machines}

\bigskip

\noindent
Dear Professor Taddeo,

\bigskip

I am writing to submit the manuscript entitled ``Geometric Ethics: Invariance, Conservation, and the Structure of Moral Evaluation'' for consideration as an original research article in \textit{Minds and Machines}.

This paper develops a mathematical framework---Geometric Ethics---that addresses a foundational question at the intersection of philosophy and AI: what mathematical structure must moral evaluation possess? The central argument is that moral evaluation is irreducibly multi-dimensional and carries geometric structure (metric, symmetry, conservation laws) that cannot be captured by the scalar representations currently dominant in AI alignment. The paper proves this claim formally via five theorems---Information Monotonicity, Structured Preservation, a Moral Noether Theorem, a $D_4$ Gauge Group characterization, and a No Escape result---and validates it empirically against 20,030 advice-column letters and 109,294 cross-lingual passages spanning 11 languages and 3,000 years.

I believe this manuscript is well suited to \textit{Minds and Machines} for three reasons:

\begin{enumerate}
    \item \textbf{Philosophy--computation intersection.} The paper bridges formal philosophy (metaethics, deontic logic, Hohfeldian analysis) and computational methodology (gauge theory, differential geometry, NLP), which is the core territory of \textit{Minds and Machines}.

    \item \textbf{Conceptual novelty with formal precision.} The framework imports gauge invariance and Noether's theorem from physics into ethics---not as metaphors, but as operational mathematical tools that generate testable predictions. The resulting ``computational objectivity'' position---that moral objectivity arises from optimization under constraint rather than transcendent grounding---offers a new contribution to metaethical debates that \textit{Minds and Machines} readers will find relevant.

    \item \textbf{Independent empirical replication.} The framework's cross-lingual invariance predictions have been independently replicated by a researcher at a separate institution using different data and methodology, confirming that moral structure is language-invariant with moral and linguistic subspaces nearly orthogonal (shared variance $< 3\%$). This strengthens the empirical grounding beyond what is typical for formal ethics contributions.
\end{enumerate}

An earlier version of this manuscript was submitted to \textit{Artificial Intelligence} (ARTINT-D-26-00335), where the Editor-in-Chief noted that the work was ``better suited to a philosophical venue.'' I have accordingly refined the presentation to foreground the philosophical contributions while retaining the formal apparatus that supports them.

This manuscript has not been published previously and is not under consideration elsewhere. The work is original, and I have read and comply with the journal's ethical guidelines.

I have suggested six reviewers with relevant expertise in the Editorial Manager system.

Thank you for considering this manuscript. I look forward to your response.

\bigskip

\noindent
Sincerely,

\bigskip

\noindent
Andrew H. Bond\\
Senior Lecturer\\
Department of Computer Engineering\\
San Jos\'{e} State University\\
\href{mailto:andrew.bond@sjsu.edu}{andrew.bond@sjsu.edu}

\end{document}
